% Chapter 4: Neymanian Repeated Sampling Inference
% A First Course in Causal Inference - Peng Ding
% Rewritten in Stein motivation-first style with textbook theorems

\chapter{Neymanian 重复抽样推断}\label{ch:4}

\section{两种推断哲学的对话}

在上一章中，我们学习了费舍尔随机化检验（FRT）。Fisher 的方法优雅而精确——它基于随机化的精确性质，在 sharp null 假设下给出精确的 p 值。

但 FRT 也有其局限：它只能检验 sharp null 假设，对于构建置信区间需要额外的推理。而且当样本量很大时，枚举所有可能的置换变得计算上不可行。

Jerzy Neyman 在 1923 年的工作中提出了一种互补的方法。这种方法基于\textbf{渐近正态分布}，可以更方便地构建置信区间和进行假设检验。Neyman 的方法不依赖于任何特定的 null 假设，而是直接对因果效应进行统计推断。

这两种方法——Fisherian 和 Neymanian——代表了统计推断的两种哲学。理解它们的联系与区别，对于真正掌握因果推断至关重要。

\userannotation{为什么需要两种方法？}{FRT 适用于小样本精确推断和假设检验，Neymanian 适用于大样本渐近推断和置信区间构建。两者互为补充，而非相互排斥。}

\section{有限总体量}

在 Neymanian 框架中，我们需要一套精确的记号来描述\textbf{有限总体量}——这些是不带 "hat" 的真实值。

考虑一个包含 $n$ 个个体的有限总体，其中 $n_1$ 个被分配到治疗组，$n_0 = n - n_1$ 个被分配到对照组。对于单元 $i = 1, \ldots, n$，我们有潜在结果 $Y_i(1)$ 和 $Y_i(0)$，以及个体效应 $\tau_i = Y_i(1) - Y_i(0)$。

\textbf{有限总体均值}：
\[
\bar{Y}(1) = n^{-1} \sum_{i=1}^n Y_i(1), \quad
\bar{Y}(0) = n^{-1} \sum_{i=1}^n Y_i(0)
\]

\textbf{有限总体方差}（使用 $n-1$ 作为除数使定理更优雅）：
\[
S^2(1) = (n-1)^{-1} \sum_{i=1}^n \{ Y_i(1) - \bar{Y}(1) \}^2, \quad
S^2(0) = (n-1)^{-1} \sum_{i=1}^n \{ Y_i(0) - \bar{Y}(0) \}^2
\]

\textbf{有限总体协方差}：
\[
S(1,0) = (n-1)^{-1} \sum_{i=1}^n \{ Y_i(1) - \bar{Y}(1) \} \{ Y_i(0) - \bar{Y}(0) \}
\]

\textbf{个体效应}：
\[
\tau = n^{-1} \sum_{i=1}^n \tau_i = \bar{Y}(1) - \bar{Y}(0)
\]
是平均因果效应的真值。

\textbf{个体效应的方差}：
\[
S^2(\tau) = (n-1)^{-1} \sum_{i=1}^n (\tau_i - \tau)^2
\]

\userannotation{为什么需要协方差？}{因为 $Y(1)$ 和 $Y(0)$ 之间可能有相关性。如果治疗组和对照组的个体是匹配的（正相关），估计的方差会减小；如果是负相关，方差会增大。}

\section{一个关键引理}

在有限总体量的世界中，有一个优雅的关系连接了协方差与各个方差：

\begin{Lemma}[Lemma 4.1 — 方差与协方差的关系]\label{lemm:4.1}
\[
2S(1,0) = S^2(1) + S^2(0) - S^2(\tau)
\]
其中 $\tau_i = Y_i(1) - Y_i(0)$ 是个体因果效应。
\end{Lemma}

这个引理的证明是直接的代数运算。这个恒等式将在后续推导中发挥关键作用。

\textbf{推论}：如果所有个体因果效应都相等（$S^2(\tau) = 0$），那么 $S(1,0) = [S^2(1) + S^2(0)]/2$，即潜在结果完全正相关。

\section{Neyman (1923) 定理}

基于观测结果，我们可以计算样本均值：
\[
\hat{\bar{Y}}(1) = n_1^{-1} \sum_{i=1}^n Z_i Y_i, \quad
\hat{\bar{Y}}(0) = n_0^{-1} \sum_{i=1}^n (1 - Z_i) Y_i
\]

样本方差：
\[
\hat{S}^2(1) = (n_1 - 1)^{-1} \sum_{i=1}^n Z_i \{ Y_i - \hat{\bar{Y}}(1) \}^2, \quad
\hat{S}^2(0) = (n_0 - 1)^{-1} \sum_{i=1}^n (1 - Z_i) \{ Y_i - \hat{\bar{Y}}(0) \}^2
\]

\textbf{关键观察}：不存在 $S(1,0)$ 和 $S^2(\tau)$ 的样本版本，因为潜在结果 $Y_i(1)$ 和 $Y_i(0)$ 从未对任何单元同时被观测到。

现在陈述 Neyman 的核心定理：

\begin{Theorem}[Neyman (1923) 定理]\label{chap4-thm:neyman1923}
在 CRE 下：
\begin{enumerate}
  \item \textbf{无偏性}：差值均值估计量 $\hat{\tau} = \hat{\bar{Y}}(1) - \hat{\bar{Y}}(0)$ 对 $\tau$ 是无偏的：
    \[
    \mathbb{E}(\hat{\tau}) = \tau
    \]

  \item \textbf{方差公式}：$\hat{\tau}$ 有方差
    \begin{align}
    \text{var}(\hat{\tau}) &= \frac{S^2(1)}{n_1} + \frac{S^2(0)}{n_0} - \frac{S^2(\tau)}{n} \label{eq:neyman-var} \\
    &= \frac{n_0}{n_1 n} S^2(1) + \frac{n_1}{n_0 n} S^2(0) + \frac{2}{n} S(1,0) \label{eq:neyman-var-equivalent}
    \end{align}

  \item \textbf{保守方差估计}：方差估计量
    \[
    \hat{V} = \frac{\hat{S}^2(1)}{n_1} + \frac{\hat{S}^2(0)}{n_0}
    \]
    对估计 $\text{var}(\hat{\tau})$ 是保守的，即
    \[
    \mathbb{E}(\hat{V}) - \text{var}(\hat{\tau}) = \frac{S^2(\tau)}{n} \geq 0
    \]
    等号当且仅当 $\tau_i = \tau$（所有单元的个体效应恒定）时成立。
\end{enumerate}
\end{Theorem}

\textbf{关于期望和方差的含义}：潜在结果都是固定数值，只有治疗指示变量 $Z_i$ 是随机的。因此，期望和方差都相对于 $Z_i$ 的随机性——即 $n_1$ 个 1 和 $n_0$ 个 0 的所有可能排列。

\userannotation{第三项为什么是减项？}{直觉上：如果个体因果效应存在变异（$S^2(\tau) > 0$），治疗组和对照组之间的差异不仅来自随机抽样，还来自因果效应本身的真实变异。这种"真实变异"不应该被算作估计误差，所以被减去。}

\section{方差公式的理解}

方差公式 \cref{eq:neyman-var} 与经典两样本问题中的方差公式不同，因为它不仅依赖于潜在结果的有限总体方差，还依赖于个体效应的有限总体方差。

\textbf{一个令人困惑的现象}：个体效应越异质，$\hat{\tau}$ 的变异性越小。为什么？

从等价格式 \cref{eq:neyman-var-equivalent} 理解：比较正相关和负相关的潜在结果情况。

假设在一次实现的实验中，我们观察到相对较大的治疗潜在结果：
\begin{itemize}
  \item 如果 $S(1,0) > 0$，则那些接受治疗的单元也有相对较大的对照潜在结果。因此，对照下的观测结果相对较小，导致较大的 $\hat{\tau}$
  \item 如果 $S(1,0) < 0$，则那些接受治疗的单元有相对较小的对照潜在结果。因此，对照下的观测结果相对较大，导致较小的 $\hat{\tau}$
\end{itemize}

总体而言，尽管 $\hat{\tau}$ 的无偏性不依赖于潜在结果之间的相关性，但在正相关情况下更可能观察到更极端的 $\hat{\tau}$。因此，当潜在结果正相关时，$\hat{\tau}$ 的方差更大。

\section{Neyman (1923) 定理的证明}\label{sec:prove-neyman1923}

\textbf{无偏性证明}：

从表示开始：
\begin{align*}
\hat{\tau} &= n_1^{-1} \sum_{i=1}^n Z_i Y_i - n_0^{-1} \sum_{i=1}^n (1 - Z_i) Y_i \\
&= n_1^{-1} \sum_{i=1}^n Z_i Y_i(1) - n_0^{-1} \sum_{i=1}^n (1 - Z_i) Y_i(0)
\end{align*}

利用期望的线性性：
\begin{align*}
\mathbb{E}(\hat{\tau}) &= \mathbb{E}\left\{ n_1^{-1} \sum_{i=1}^n Z_i Y_i(1) - n_0^{-1} \sum_{i=1}^n (1 - Z_i) Y_i(0) \right\} \\
&= n_1^{-1} \sum_{i=1}^n \mathbb{E}(Z_i) Y_i(1) - n_0^{-1} \sum_{i=1}^n \mathbb{E}(1 - Z_i) Y_i(0) \\
&= n_1^{-1} \sum_{i=1}^n \frac{n_1}{n} Y_i(1) - n_0^{-1} \sum_{i=1}^n \frac{n_0}{n} Y_i(0) \\
&= n^{-1} \sum_{i=1}^n Y_i(1) - n^{-1} \sum_{i=1}^n Y_i(0) = \tau
\end{align*}

\textbf{方差证明}：

我们可以进一步将 $\hat{\tau}$ 写为：
\[
\hat{\tau} = \sum_{i=1}^n Z_i \left\{ \frac{Y_i(1)}{n_1} + \frac{Y_i(0)}{n_0} \right\} - n_0^{-1} \sum_{i=1}^n Y_i(0)
\]

利用简单随机抽样的引理，得到：
\begin{align*}
\text{var}(\hat{\tau}) &= \frac{n_1 n_0}{n(n-1)} \sum_{i=1}^n \left\{ \frac{Y_i(1)}{n_1} + \frac{Y_i(0)}{n_0} - \frac{\bar{Y}(1)}{n_1} - \frac{\bar{Y}(0)}{n_0} \right\}^2 \\
&= \frac{n_0}{n_1 n} S^2(1) + \frac{n_1}{n_0 n} S^2(0) + \frac{2}{n} S(1,0)
\end{align*}

利用 \cref{lemm:4.1}，我们也可以将方差写为：
\[
\text{var}(\hat{\tau}) = \frac{S^2(1)}{n_1} + \frac{S^2(0)}{n_0} - \frac{S^2(\tau)}{n}
\]

\textbf{方差估计的期望}：

因为治疗组是从 $n$ 个单元中简单随机抽取的 $n_1$ 个样本，简单随机抽样引理确保：
\[
\mathbb{E}\{\hat{S}^2(1)\} = S^2(1), \quad \mathbb{E}\{\hat{S}^2(0)\} = S^2(0)
\]

因此，$\hat{V}$ 对 \cref{eq:neyman-var} 中前两项是无偏的。

\section{渐近正态性与置信区间}

\cite{li2017} 进一步证明了基于有限总体 CLT 的 $\hat{\tau}$ 的渐近正态性：

\begin{Theorem}[CLT for 差值均值]\label{chap4-thm:clt}
令 $n \to \infty$ 和 $n_1 \to \infty$。如果 $n_1/n$ 有 $(0, 1)$ 中的极限值，$\{S^2(1), S^2(0), S(1,0)\}$ 有极限值，且
\[
\max_{1 \leq i \leq n} \{Y_i(1) - \bar{Y}(1)\}^2 / n \to 0, \quad
\max_{1 \leq i \leq n} \{Y_i(0) - \bar{Y}(0)\}^2 / n \to 0
\]
则
\[
\frac{\hat{\tau} - \tau}{\sqrt{\text{var}(\hat{\tau})}} \xrightarrow{d} N(0, 1)
\]
且
\[
\hat{S}^2(1) \to S^2(1), \quad \hat{S}^2(0) \to S^2(0) \quad \text{（依概率）}
\]
\end{Theorem}

这个定理确保当样本量足够大时，$\hat{\tau}$ 的抽样分布可以用正态分布近似。它还确保样本方差是总体方差的一致估计。

\textbf{置信区间}：这为 $\tau$ 的保守大样本置信区间提供了正当理由：
\[
\hat{\tau} \pm z_{1-\alpha/2} \sqrt{\hat{V}}
\]
其中 $z_{1-\alpha/2}$ 是标准正态分布的 $1 - \alpha/2$ 分位数。

\textbf{与弱零假设的對偶性}：通过置信区间与假设检验的对偶性，这也提供了检验
\[
H_0: \tau = 0
\]
（称为弱零假设）的基础。

\userannotation{为什么 Neymanian 方法只给出"弱零假设"的检验？}{因为在 Neymanian 框架中，我们无法观测到 $S^2(\tau)$，所以无法精确计算方差。置信区间和检验都是渐近有效的，只在样本量大时可靠。}

\section{Neymanian 推断与 OLS 的关系}

实践中，研究者经常使用基于回归的推断来估计平均因果效应。标准方法是运行 OLS 回归：
\[
(\hat{\alpha}, \hat{\beta}) = \arg\min_{(a,b)} \sum_{i=1}^n (Y_i - a - b Z_i)^2
\]
并使用 treatment 系数的估计量 $\hat{\beta}$ 作为平均因果效应的估计量。

\textbf{关键发现}：系数 $\hat{\beta}$ 恰好等于差值均值：
\[
\hat{\beta} = \hat{\tau}
\]

但是，OLS 的通常方差估计量与 $\hat{V}$ 不同。幸运的是，Eicker-Huber-White (EHW) 稳健方差估计量接近 $\hat{V}$，HC2 变体恰好等于 $\hat{V}$。

在 LaLonde 数据上的应用：
\begin{verbatim}
> tauhat = mean(y[z==1]) - mean(y[z==0])
> vhat   = var(y[z==1])/n1 + var(y[z==0])/n0
> sehat  = sqrt(vhat)
> tauhat
[1] 1794.343
> sehat
[1] 670.9967
\end{verbatim}

OLS 的标准误（通常）太小：
\begin{verbatim}
> olsfit = lm(y ~ z)
> summary(olsfit)$coef[2, 1:2]
  Estimate Std. Error
 1794.3431   632.8536
\end{verbatim}

使用 EHW 稳健标准误得到正确结果：
\begin{verbatim}
> sqrt(hccm(olsfit, type = "hc2")[2, 2])
[1] 670.9967
\end{verbatim}

\section{重尾结局与正态近似的失效}

\cref{chap4-thm:clt} 中的 CLT 在某些正则条件下成立。当潜在结局为重尾分布时，这些条件可能被违反。

当对照潜在结果被污染时（例如 Cauchy 分布），正态近似效果很差。模拟研究表明，在重尾情况下，直方图和正态近似之间存在显著差异。

\section{本章小结}

本章建立了 Neymanian 推断的基础：

\begin{enumerate}
  \item \textbf{有限总体量}：不带 hat 的真实值
    \[
    \bar{Y}(1), \bar{Y}(0), S^2(1), S^2(0), S(1,0), S^2(\tau)
    \]
  \item \textbf{Lemma 4.1}：$2S(1,0) = S^2(1) + S^2(0) - S^2(\tau)$
  \item \textbf{Neyman (1923) 定理}：
    \begin{itemize}
      \item $\mathbb{E}(\hat{\tau}) = \tau$（无偏）
      \item $\text{var}(\hat{\tau}) = \frac{S^2(1)}{n_1} + \frac{S^2(0)}{n_0} - \frac{S^2(\tau)}{n}$
      \item $\mathbb{E}(\hat{V}) = \text{var}(\hat{\tau}) + \frac{S^2(\tau)}{n}$（保守）
    \end{itemize}
  \item \textbf{方差估计}：$\hat{V} = \frac{\hat{S}^2(1)}{n_1} + \frac{\hat{S}^2(0)}{n_0}$
  \item \textbf{置信区间}：$\hat{\tau} \pm z_{1-\alpha/2} \sqrt{\hat{V}}$
  \item \textbf{CLT}：$\frac{\hat{\tau} - \tau}{\sqrt{\text{var}(\hat{\tau})}} \xrightarrow{d} N(0, 1)$
  \item \textbf{OLS 关系}：$\hat{\beta} = \hat{\tau}$，但需使用 EHW 稳健标准误
\end{enumerate}

下一章我们将学习如何通过分层（stratification）来提高估计的精度——这是一种利用协变量信息来改进推断的方法。

\section{习题}

\begin{Exercise}{\ref{ex:4.1} 引理的证明}\label{ex:4.1}
证明 Lemma \ref{lemm:4.1}：$2S(1,0) = S^2(1) + S^2(0) - S^2(\tau)$。
\end{Exercise}

\begin{Exercise}{\ref{ex:4.2} Neyman (1923) 定理的另一种证明}\label{ex:4.2}
在 CRE 下，计算：
\[
\text{var}\{\hat{\bar{Y}}(1)\}, \quad \text{var}\{\hat{\bar{Y}}(0)\}, \quad \text{cov}\{\hat{\bar{Y}}(1), \hat{\bar{Y}}(0)\}
\]
并利用这些公式计算 $\text{var}(\hat{\tau})$。
\end{Exercise}

\begin{Exercise}{\ref{ex:4.3} Neymanian 推断与 OLS}\label{ex:4.3}
证明 $\hat{\beta} = \hat{\tau}$（即 OLS 系数等于差值均值）。证明 EHW 稳健方差估计量的 HC2 变体恰好等于 $\hat{V}$。
\end{Exercise}

\begin{Exercise}{\ref{ex:4.4} 治疗效应的异质性}\label{ex:4.4}
证明 $S^2(\tau) = 0$ 蕴含 $S^2(1) = S^2(0)$。给出反例说明 $S^2(1) = S^2(0)$ 但 $S^2(\tau) \neq 0$。

证明 $S^2(1) < S^2(0)$ 蕴含 $S(Y(0), \tau) < 0$。给出反例说明 $S^2(1) > S^2(0)$ 但 $S(Y(0), \tau) < 0$。
\end{Exercise}

\begin{Exercise}{\ref{ex:4.5} 方差公式的改进上界}\label{ex:4.5}
证明方差公式的改进上界：
\[
\text{var}(\hat{\tau}) \leq \frac{1}{n}\left\{ \sqrt{\frac{n_0}{n_1}}S(1) + \sqrt{\frac{n_1}{n_0}}S(0) \right\}^2
\]
并讨论等号成立的条件。
\end{Exercise}

\begin{Exercise}{\ref{ex:4.6} Neyman (1923) 的向量版本}\label{ex:4.6}
将 Neyman (1923) 的结果推广到多结局情形。考虑向量结局 $\boldsymbol{V} \in \mathbb{R}^K$ 的平均因果效应：
\[
\tau_{\boldsymbol{V}} = \frac{1}{n}\sum_{i=1}^n \{\boldsymbol{V}_i(1) - \boldsymbol{V}_i(0)\}
\]

证明 $\hat{\tau}_{\boldsymbol{V}}$ 对 $\tau_{\boldsymbol{V}}$ 是无偏的，并求出其协方差矩阵。
\end{Exercise}

\begin{Exercise}{\ref{ex:4.7} BRE 中的推断}\label{ex:4.7}
回忆在 \cref{def:cre} 中定义的完全随机实验。在这种设定下，推导平均因果效应的方差公式，并比较其与 CRE 中方差的差异。
\end{Exercise}

\begin{Exercise}{\ref{ex:4.8} 推荐阅读}\label{ex:4.8}
\begingroup
\parskip=0pt
\textbf{推荐阅读}：\cite{ding2017} 重新审视了因果推断中的一些经典问题。\cite{imbens2015} 是关于因果推断的权威综述。

\textbf{历史注记}：Neyman 1923 年的原始论文\cite{neyman1923}奠定了有限总体推断的基础。当代统计学家将其视为因果推断范式的里程碑。
\endgroup
\end{Exercise}

% === 用户问答记录 ===
% (后续通过 interactive-qa 更新)
